\documentclass[runningheads]{llncs}

\usepackage[T1]{fontenc}
\usepackage{graphicx}
\usepackage{amsmath,amssymb}
\usepackage{booktabs}
\usepackage{multirow}
\usepackage{array}
\usepackage{xcolor}
\usepackage{tikz}
\usetikzlibrary{arrows.meta,positioning,fit,calc,backgrounds}
\usepackage{subcaption}
\usepackage[hidelinks]{hyperref}
\usepackage[capitalise,noabbrev]{cleveref}
\usepackage{listings}
\usepackage{placeins}

\renewcommand\UrlFont{\color{blue}\rmfamily}
\lstset{basicstyle=\ttfamily\scriptsize,columns=fullflexible,frame=single,
        framerule=0.3pt,rulecolor=\color{gray},breaklines=true,xleftmargin=2pt}

\makeatletter\newcommand{\tabinput}[1]{\@@input #1 }\makeatother
\newcommand{\R}{\mathbb{R}}
\newcommand{\E}{\mathbb{E}}
\newcommand{\sg}{\operatorname{sg}}
\newcommand{\gp}[1]{\textsc{#1}}

% palette used in every figure (validated categorical slots 1-5)
\definecolor{cNat}{HTML}{2A78D6}
\definecolor{cFsqF}{HTML}{EB6834}
\definecolor{cFsqJ}{HTML}{1BAF7A}
\definecolor{cJep}{HTML}{EDA100}
\definecolor{cGray}{HTML}{8A8A8A}

\setlength{\textfloatsep}{9pt plus 2pt minus 3pt}
\setlength{\floatsep}{7pt plus 2pt minus 2pt}
\setlength{\intextsep}{7pt plus 2pt minus 2pt}

\begin{document}

\title{Towards a Universal Feature Extractor for Atari Games:\\
Reconstruction, Finite Scalar Quantisation and\\ Joint-Embedding Prediction}
\titlerunning{Towards a Universal Feature Extractor for Atari Games}

\author{Sajid Ali Sadruddin}
\authorrunning{Sajid Ali Sadruddin}
\institute{TU Dresden, Germany\\[2pt]
Research Project, supervised by Oleh Shkalikov}

\maketitle

\begin{abstract}
Reinforcement learning (RL) agents for Atari games usually learn a new convolutional
feature extractor for every single game, even though all games give the agent the same kind
of input: four stacked $84\times84$ grey-scale frames. In this project, I investigated whether
one encoder, pre-trained once on expert gameplay from eleven games, can be reused as a
universal feature extractor for Proximal Policy Optimisation (PPO). I compared autoencoders
with vector-quantised and Finite Scalar Quantisation (FSQ) bottlenecks against an action- and
game-conditioned Joint-Embedding Predictive Architecture (V-JEPA) world model. Each encoder was
used either frozen or fine-tuned together with the policy for 10\,M environment steps on
eleven seen and four unseen games. The fine-tuned FSQ-VAE reached a higher median
human-normalised score than a NatureCNN trained from scratch (1.07 vs.\ 0.74), beat it on four
of eleven games, and on four games reached 80\% of its final level 1.3--2.7\,M steps earlier.
Kept frozen, the same encoder reached less than a third of that median. The V-JEPA learned a
healthy latent space and real dynamics on the seen games, but PPO could hardly use its
features (median 0.05). My analysis shows that the model mostly ignored the action, produced
very similar latents for consecutive frames and specialised to the pre-training games. Overall,
a shared pre-trained encoder works well as a starting point, but not yet as a frozen
replacement for the NatureCNN.
\end{abstract}

% =========================================================================
\section{Introduction}
\label{sec:intro}

Atari 2600 games, played through the Arcade Learning Environment~\cite{machado2018revisiting},
are one of the most common benchmarks for learning control from pixels. Since
DQN~\cite{mnih2015human}, the standard agent stacks the last four pre-processed frames and
passes them through a small convolutional network, the NatureCNN, followed by a policy and a
value head. All of this is trained from the game reward alone, and it is done again from
scratch for every game. This is wasteful, because the games share the same input format and a similar
visual vocabulary of sprites, balls, walls and scrolling backgrounds (\cref{fig:games}). A
feature extractor that is trained once and then reused could save interactions and
computation, transfer to new games, and give one latent space for all games.

This project was proposed by Oleh Shkalikov. The question is whether such a
\emph{universal feature extractor} can be trained for Atari, and I looked at it from two sides.
The first is the pre-training objective: I used quantised
autoencoders~\cite{vandenoord2017vqvae,mentzer2024fsq}, which learn by reconstructing the
frames, and a joint-embedding predictive architecture
(JEPA)~\cite{lecun2022path,assran2023ijepa,bardes2024vjepa}, which learns by predicting the
next state in latent space. The second is how the extractor is used: kept frozen, or fine-tuned
together with the policy. Specifically, I wanted to answer four questions: which objective
gives the most useful features for PPO (\textbf{RQ1}), whether freezing or fine-tuning is
better (\textbf{RQ2}), whether the features generalise to games that were not seen during
pre-training (\textbf{RQ3}), and what the approach costs and how interpretable it is
(\textbf{RQ4}).

% =========================================================================
\section{Related Work}
\label{sec:related}

PPO~\cite{schulman2017ppo} with the RL-Zoo hyper-parameters~\cite{raffin2020zoo} is the
standard on-policy baseline for Atari. Many multi-game agents share the whole policy across
games, while I only share the perception part, so that the effect of the representation can
be measured on its own. Discrete latent spaces work well on Atari
(DreamerV2~\cite{hafner2021dreamerv2}, IRIS~\cite{micheli2023iris}), and self-predictive methods
such as ATC~\cite{stooke2021atc} and SPR~\cite{schwarzer2021spr} avoid reconstructing pixels.
ATC even showed that an encoder trained without rewards can be frozen for some Atari games. In
robotics, pre-trained encoders have been reported both as very effective~\cite{parisi2022pvr}
and as no better than an encoder trained from scratch~\cite{hansen2022pretrained}, a tension I
also ran into. For the JEPA, I borrowed the learnable task embeddings of the TD-MPC2 world
model~\cite{hansen2024tdmpc2}.

% =========================================================================
\section{Prerequisites}
\label{sec:prereq}

\paragraph{Problem Formulation.} In every game $g$, the agent sees a frame stack
$x_t\in[0,1]^{4\times84\times84}$ and can choose from its own set of 4 to 18 actions. A feature
extractor $f_\theta$ maps $x_t$ to a feature vector, on which PPO trains an actor and a critic.
The goal is one set of weights $\theta^\star$, learned without rewards on eleven games, that
works for PPO on any game. I compare a NatureCNN trained from scratch (the baseline) with a
pre-trained encoder that is either kept \emph{frozen} or \emph{fine-tuned} with the policy.

\paragraph{Dataset.} Random play rarely gets past the first screen of a game, so I collected
the data with pre-trained agents from the HuggingFace Hub (\cref{fig:data}a). Their actions were
sampled from their policies, so the data also contains some sub-optimal behaviour. I built
two datasets. The \emph{frame dataset} for the
autoencoders contains 50\,k frame stacks per game, which gives 550\,k stacks from the eleven
pre-training games and 200\,k from the four held-out games. The \emph{transition dataset} for
the world model contains tuples $(s_t,a_t,s_{t+1},d_t)$: 250\,k per pre-training game (2.75\,M in
total), a validation split of 10\,k per game and 50\,k per held-out game, each collected with a
different seed. When an episode ends, I store the real last frame instead of the reset frame.
Consecutive stacks share three of four frames, so the next state is already 75\% identical to
the current one, and transitions at the end of an episode are masked out (\cref{fig:data}b).

\begin{figure}[t]
\centering
\input{figures/data_pipeline}
\caption{(a)~Data collection with pre-trained expert agents and the two resulting corpora.
(b)~Processing at training time: frame-stack overlap, masking of episode-ending
transitions, mapping of game-specific action indices to the canonical ALE-18 set, and the
motion mask of the FSQ-VAE loss.}
\label{fig:data}
\end{figure}

% =========================================================================
\section{The Games}
\label{sec:games}

\Cref{fig:games} shows all 15 games as the encoder sees them. Eleven are used for pre-training
and four are held out for testing generalisation. Every game produces a tensor of the same
shape, but its content differs a lot (\cref{tab:games}), and one shared encoder has to handle
all of these differences at once.

\begin{figure}[t]
\centering
\includegraphics[width=\textwidth]{figures/games_grid.pdf}
\caption{All 15 games as the encoder sees them: the newest frame $t$ of an expert frame stack
(left of each pair) and $|x_{t+1}-x_t|$, the change to the next frame (right, inverted, per-image
contrast). $^\dagger$\,held-out games. Frames from the data specimen sheet of the transition
corpus.}
\label{fig:games}
\end{figure}

\begin{table}[t]
\centering
\caption{How the games differ. $|\mathcal{A}|$: minimal action set; pixel mean (0--255) and
black \% of the frames; moving \%: pixels that change within a stack ($\delta=0.05$); steps per
terminal (the Atari wrappers also end an episode on a life loss); top action: share of the
expert's most frequent action. 50\,k transitions per game; $^\dagger$\,held-out.}
\label{tab:games}
\scriptsize\setlength{\tabcolsep}{2.2pt}
\begin{tabular}{@{}llccrrrrr@{}}
\toprule
Game & Type & $|\mathcal{A}|$ & Expert & Pix. & Black & Mov. & Steps/ & Top \\
 & & & & mean & \% & \% & term. & act.\ \% \\
\midrule
\tabinput{tables/games_rows.tex}
\bottomrule
\end{tabular}
\end{table}

\paragraph{Appearance.} The average brightness goes from 5 (DemonAttack) to 105 (Pong), and
the share of completely black pixels from 1\% (Pong) to 95\% (Asteroids). Shooting games are
mostly empty space, while sports and scrolling games fill the whole screen.

\paragraph{Motion.} In most games, less than 4\% of the pixels change within a frame stack
(Pong 1.7\%, Breakout 1.4\%, IceHockey 1.1\%). Only the scrolling and pseudo-3D games move large
parts of the screen (Riverraid 10\%, BeamRider 34\%). So the objects that matter for playing
are usually a few small sprites on a static background, which a loss that treats every pixel
equally can almost ignore.

\paragraph{Actions and episode length.} The number of actions ranges from 4 to 18, and the same
action index means a different joystick command in different games. A shared action embedding
therefore needs a common action set (\cref{sec:models}). Episodes last
from about 120 steps (Atlantis) to 12\,500 steps (Enduro). The expert data is also not equally
varied: the expert's most frequent action covers 8\% of the steps in Seaquest, but 67\% in
RoadRunner and 99\% in Asteroids. For these games, the data rarely shows what would have
happened after a different action.

% =========================================================================
\section{Models}
\label{sec:models}

All encoders take the same input $x\in[0,1]^{4\times84\times84}$. After pre-training I throw
away the decoder or predictor and only keep the encoder, whose output is passed to the actor
and critic of PPO. The baseline is the NatureCNN with a 512-dimensional feature, trained from
scratch.

\subsection{Autoencoders}

I started with a Gaussian VAE~\cite{kingma2014vae} with the NatureCNN layers as encoder,
but only trained it for one epoch: its $7\times7$ latent grid makes each cell responsible for a
$12\times12$ pixel region, much larger than a ball or a bullet. I therefore moved to spatial,
discrete bottlenecks.

\paragraph{Residual VQ-VAE.} This model keeps a finer $21\times21$ latent grid, using two
stride-2 convolutions and two residual blocks. Each 64-dimensional latent cell is replaced by
the closest of $K=512$ codebook vectors, and the gradient is passed through with the
straight-through trick. The codebook is updated with exponential moving averages
($\gamma=0.99$), the commitment weight is $\beta=0.25$, and codes that are no longer used are
re-initialised. When I looked at the reconstructions, small moving objects were often missing.
To fix this, I trained a second version in which the error of moving pixels,
$m=\mathbf{1}[\max_j|x^{(j+1)}-x^{(j)}|>0.05]$, counts ten times more.

\paragraph{FSQ-VAE.} Finite Scalar Quantisation~\cite{mentzer2024fsq} removes the codebook
completely. A three-layer convolutional encoder produces four channels, and each channel is
rounded to a fixed number of levels, $L=[8,5,5,3]$, which gives 600 possible values per latent
cell. There is no codebook, no commitment loss and no unused codes. My first version used
transposed convolutions in the decoder, which produced checkerboard patterns on the flat Atari
backgrounds, so I replaced them with upsampling followed by normal convolutions. Instead of the
fixed factor of ten, the final loss treats moving and static pixels separately:
\begin{equation}
  \mathcal{L}_{\text{bal}} = w\,\tfrac{\sum m\odot(\hat x-x)^2}{\sum m}
  +(1-w)\,\tfrac{\sum (1-m)\odot(\hat x-x)^2}{\sum (1-m)},\qquad w=0.85.
\end{equation}
This way, the moving pixels always get 85\% of the loss, no matter how few of them there are.
For PPO, the feature is the flattened $4\times21\times21$ quantised grid (1\,764 values).

\subsection{Action- and Game-Conditioned V-JEPA}

A JEPA does not reconstruct pixels. Instead, it learns to predict the \emph{representation}
of the next observation~\cite{lecun2022path}. If the prediction also receives the action, the
model becomes a latent world model, similar to SPR~\cite{schwarzer2021spr}. \Cref{fig:jepa}
shows my final model. A $14\times14$ patch embedding turns a frame stack into 36 tokens of
size 256. The context encoder $E_\theta$, a six-layer transformer, encodes $s_t$, and the
target encoder $\bar E$, an exponential moving average (EMA) of $E_\theta$ with $\tau=0.996$
that receives no gradients~\cite{assran2023ijepa}, encodes $s_{t+1}$. A three-layer predictor
$P_\psi$ gets the context tokens together with an action token and a game token, and has to
predict the target tokens:
\begin{equation}
  \mathcal{L}_{\text{JEPA}}=\frac{\sum_i m_i\,\big\|P_\psi\big(E_\theta(s_t^i,e^i),a_t^i,e^i\big)
  -\operatorname{sg}\big(\bar E(s_{t+1}^i,e^i)\big)\big\|_2^2}{\sum_i m_i},
\end{equation}
where $m_i=0$ for transitions that end an episode.

My first version trained well according to its loss curve, but a detailed check of the
pipeline showed that action indices meant different commands in different games, the large
dataset was loaded with the wrong axis order, and the validation set only contained one game.
I rebuilt the model (v2) following TD-MPC2~\cite{hansen2024tdmpc2}. First, all actions are mapped to the common set of 18 ALE
actions. Second, each game gets a learnable embedding $e_g$ with $\|e_g\|\le1$, which is added
as a token to both encoders and the predictor. With probability 0.25 it is replaced by a shared
UNKNOWN embedding, so that the model also learns a general ``Atari'' prior for unseen games.
Third, all latents are normalised with SimNorm, which splits them into 32 groups of $V=8$ values
and applies a softmax to each group. This keeps them bounded and makes collapse harder. Fourth,
an inverse-dynamics head~\cite{pathak2017icm} predicts the action $a_t$ from $z_t$,
$\bar z_{t+1}$ and $e_g$, only among the actions that are valid in the game. The total loss is
$\mathcal{L}=\mathcal{L}_{\text{JEPA}}+\lambda\mathcal{L}_{\text{inv}}$ with $\lambda=0.01$.

\paragraph{Training Step.} The target branch is computed without gradients, both losses are
averaged over the same unmasked samples, the gradients are clipped to a norm of 1.0, AdamW
updates the trainable parts, and only then is the target encoder moved slightly towards the
context encoder (\cref{tab:grad}). The target is therefore a slowly moving goal that the context
encoder cannot change directly, which prevents collapse to a constant output. Note also that
the inverse loss trains the encoder but \emph{not} the predictor, because its second input is
the target latent and not the prediction.

\paragraph{Using the Encoder in PPO.} For PPO, the context encoder receives the game's own
embedding (or UNKNOWN for unseen games) and produces $36\times256$ tokens. In the main runs I
flatten them into a 9\,216-dimensional feature, and in a second set of runs I average them into
a 256-dimensional feature.

\begin{table}[t]
\centering
\caption{Gradient routing in the V-JEPA (12.48\,M parameters, 7.53\,M trainable).}
\label{tab:grad}
\footnotesize\setlength{\tabcolsep}{4pt}
\begin{tabular}{@{}lrccl@{}}
\toprule
Parameter group & Params & $\nabla\mathcal{L}_{\text{JEPA}}$ & $\nabla\mathcal{L}_{\text{inv}}$ & Update \\
\midrule
Context encoder (patchify, 6 blocks, LN) & 4.95\,M & yes & yes & AdamW \\
Predictor (action emb., 3 blocks, LN) & 2.37\,M & yes & no & AdamW \\
Inverse head ($768\to256\to18$) & 0.20\,M & no & yes & AdamW \\
Game embeddings ($16\times256$) & 0.004\,M & yes & yes & AdamW \\
Target encoder (copy of context enc.) & 4.95\,M & no & no & EMA only \\
\bottomrule
\end{tabular}
\end{table}

\begin{figure}[t]
\centering
\input{figures/jepa_arch}
\caption{Environment- and action-conditioned V-JEPA (v2), as implemented. The target
encoder is an EMA copy of the context encoder; only the context encoder is kept for PPO.}
\label{fig:jepa}
\end{figure}

% =========================================================================
\section{Methodology}
\label{sec:method}

\subsection{Training Configurations}

\begin{table}[t]
\centering
\caption{Training configurations of the four encoders used with PPO (``this work'') and
the setting of the reference paper each follows. ``--'': not specified, or not logged.}
\label{tab:pretrain}
\scriptsize\setlength{\tabcolsep}{2.5pt}
\newcolumntype{Y}[1]{>{\raggedright\arraybackslash}p{#1}}
\begin{tabular}{@{}Y{1.6cm}Y{2.25cm}Y{2.2cm}Y{2.95cm}Y{2.25cm}@{}}
\toprule
 & Residual VQ-VAE & FSQ-VAE & V-JEPA v2 & NatureCNN \\
\midrule
\multicolumn{5}{@{}l}{\textit{This work}} \\
Objective & MSE$^{a}$ + VQ + commit. & $\mathcal{L}_{\text{bal}}$, $w=0.85$ &
  $\mathcal{L}_{\text{JEPA}}+0.01\,\mathcal{L}_{\text{inv}}$ & PPO only (no pre-training) \\
Data & 550\,k stacks & 550\,k stacks & 2.75\,M transitions & online, $10^7$ steps/game \\
Length & 15 ep.\ (129\,k steps) & 125 ep.\ (537\,k steps) & 92 ep.\ (21.2\,k steps) & 9.8\,k rollouts \\
Batch & 64 & 128 & 512 & $8\times128$; minibatch 256 \\
Optimiser, LR & $3{\cdot}10^{-4}$ & AdamW, $3{\cdot}10^{-4}$ & AdamW, $3{\cdot}10^{-4}$ & Adam, $2.5{\cdot}10^{-4}\to0$ \\
Weight decay & -- & $10^{-4}$ & $10^{-5}$ (not LN/emb.) & 0 \\
Schedule & -- & 100 ep.\ const.\ + cosine & warm-up + cosine & linear to 0 \\
Bottleneck / target & $K{=}512$, $d{=}64$, $\beta{=}0.25$, EMA 0.99 & FSQ $L{=}[8,5,5,3]$ (600) &
  EMA $\tau{=}0.996$ (fixed); SimNorm $V{=}8$; patch 14 & -- \\
Parameters & 0.70\,M & 0.84\,M & 12.5\,M & 1.68\,M \\
Hardware, time & Tesla T4, 3.0\,h & H100, 1.9\,h & H100, 15.6\,h & H100, 3.0\,h/game \\
\midrule
\multicolumn{5}{@{}l}{\textit{Reference setting}} \\
Source & VQ-VAE~\cite{vandenoord2017vqvae} & FSQ~\cite{mentzer2024fsq} &
  I-JEPA~\cite{assran2023ijepa} / V-JEPA~\cite{bardes2024vjepa} / TD-MPC2~\cite{hansen2024tdmpc2} &
  PPO~\cite{schulman2017ppo}, RL-Zoo~\cite{raffin2020zoo} \\
Length & 250\,k steps & -- & 300 ep.\ / 90\,k it. / -- & $10^7$ steps (40\,M frames) \\
Batch & 128 & -- & 2\,048 / 3\,072 / 256 & $8\times128$; minibatch $32\times8$ \\
Optimiser, LR & Adam, $2{\cdot}10^{-4}$ & -- & AdamW, $10^{-3}$ / $6.25{\cdot}10^{-4}$ / Adam, $3{\cdot}10^{-4}$ &
  Adam, $2.5{\cdot}10^{-4}\cdot\alpha$ \\
Weight decay & -- & -- & $0.04\to0.4$ (I-, V-JEPA) & -- \\
Bottleneck / target & $K{=}512$, $\beta{=}0.25$, EMA 0.99 & $[8,6,5]$ ($2^8$), $[8,5,5,5]$ ($2^{10}$) &
  EMA $0.996\to1$ / $0.998\to1$; SimNorm $V{=}8$; patch 14 / 16 & 3 epochs, VF coef.\ 1 (Zoo: 4, 0.5) \\
\bottomrule
\end{tabular}\\[2pt]
{\scriptsize $^{a}$\,plain MSE, or with moving pixels weighted $10\times$ (sprite-weighted variant).}
\end{table}

\Cref{tab:pretrain} shows how I trained each encoder, next to the setting of the paper it
follows. The architectures follow the references. The differences come from scale: the JEPA
papers train very large vision transformers on huge datasets, while my model has 12.5\,M
parameters and was trained on one GPU, so I used a smaller batch, a lower learning rate, less
weight decay and a fixed EMA momentum. For the V-JEPA, I used the checkpoint with the lowest
validation loss (epoch 90 of 92), and a small sweep over $\lambda\in\{0.003,0.01,0.03,0.1\}$
confirmed $\lambda=0.01$.

\subsection{PPO Training and Evaluation}

All PPO agents use the standard RL-Zoo Atari settings (last column of \cref{tab:pretrain},
with $\gamma=0.99$, $\lambda_{\text{GAE}}=0.95$, clip range 0.1 and entropy coefficient 0.01)
and the same 512-unit actor and critic heads. The FSQ agents (15 games) were trained with a
CleanRL-style PPO loop, and the NatureCNN and V-JEPA agents (11 seen games) with
Stable-Baselines3. To measure performance, I use the mean training return over the last 10\%
of training (steps 9--10\,M). To compare games with very different score ranges, I convert it to
the human-normalised score HNS$=(\text{score}-\text{random})/(\text{human}-\text{random})$ with
the reference scores from Badia et al.~\cite{badia2020agent57}, and I mainly report the median
over games. All numbers are computed from the complete logs in Weights \& Biases. Because every
configuration was run with only one seed, small differences should not be over-interpreted~\cite{agarwal2021precipice}.

% =========================================================================
\section{Results}
\label{sec:results}

\subsection{Pre-Training}
\label{sec:pretrain-res}

\begin{figure}[t]
\centering
\includegraphics[width=\textwidth]{figures/pretrain_curves.pdf}
\caption{Pre-training curves. (a)~FSQ-VAE reconstruction MSE on the held-out games per
epoch. (b--d)~V-JEPA: latent loss on seen (true game id) and held-out games (UNKNOWN),
effective rank of the context latents, and inverse-dynamics accuracy.}
\label{fig:pretrain}
\end{figure}

The pre-training curves (\cref{fig:pretrain}) showed two things I did not expect. First, the
error on the held-out games was lowest very early. The FSQ-VAE had its lowest held-out error
after only three epochs and ended 21\% higher after 125 epochs. For the V-JEPA, the held-out
loss ended three times higher than on the seen games, and the held-out inverse accuracy dropped
below chance. Both models therefore specialised to the pre-training games. Second, a rising
loss was not a bad sign for the V-JEPA. Its validation loss was lowest in the first epoch and
then went up, while the effective rank of its latents grew from 1.4 to 57 (out of 256) and the
inverse accuracy from 23\% to 58\% (chance: 13\%). The low loss at the start simply came from an
almost collapsed latent space, where prediction is easy.

For reconstruction, I evaluated the autoencoders on 256 stacks per game. On the seen games the
FSQ-VAE had the lowest error ($0.37\cdot10^{-3}$) and three times less error on moving pixels than
the plain VQ-VAE, so the balanced loss and the new decoder did help with small sprites.
Weighting the sprites in the VQ-VAE reduced the moving-pixel error by 27\%, but increased the
overall error by 19\%. On the held-out games, all errors were four to eleven times higher and the
FSQ-VAE was the worst, so none of the autoencoders generalised to new games at the pixel level.

\subsection{Downstream Reinforcement Learning}
\label{sec:rl}

\begin{table}[t]
\centering
\caption{Mean episode return over the last 10\% of 10\,M training steps (single seed); best
per game in bold. $^\dagger$\,held-out games (FSQ runs only). V-JEPA with the flattened
feature; the pooled variant is in \cref{tab:jepa-abl}.}
\label{tab:scores}
\footnotesize\setlength{\tabcolsep}{3.8pt}
\begin{tabular}{@{}lccccc@{}}
\toprule
 & NatureCNN & \multicolumn{2}{c}{FSQ-VAE} & \multicolumn{2}{c}{V-JEPA} \\
\cmidrule(lr){3-4}\cmidrule(lr){5-6}
Game & scratch & frozen & joint & frozen & joint \\
\midrule
\tabinput{tables/scores_rows.tex}
\bottomrule
\end{tabular}
\end{table}

\begin{figure}[t]
\centering
\includegraphics[width=\textwidth]{figures/learning_curves.pdf}
\caption{Learning curves over 10\,M agent steps (100 bins; FSQ curves smoothed over five).
Solid: trained or fine-tuned encoders; dashed: frozen encoders.}
\label{fig:curves}
\end{figure}

\paragraph{FSQ-VAE, Fine-Tuned.} Fine-tuning the pre-trained FSQ encoder gave the best agent
overall (\cref{tab:scores,fig:curves}). It matched or beat the NatureCNN on Breakout, Enduro,
Qbert and Riverraid, was within 1\% on Pong and Seaquest, and reached a higher median HNS (1.07
vs.\ 0.74). The NatureCNN stayed ahead on DemonAttack, RoadRunner, BeamRider and SpaceInvaders,
which is why its mean HNS is higher (2.51 vs.\ 2.02). Where the pre-training helped, it helped a
lot (\cref{tab:sampleeff}): on Enduro, MsPacman, Qbert and Riverraid, the fine-tuned agent
reached 80\% of the NatureCNN's final level 1.3--2.7\,M steps (18--39\%) earlier, and at the
50\% level it was faster on five games.

\begin{table}[t]
\centering
\caption{Sample efficiency: million steps until the learning curve first covers 50\% and 80\%
of the way from NatureCNN's initial to its final return (``--'': never). FSQ-J/F: FSQ-VAE
joint/frozen. Fastest in bold.}
\label{tab:sampleeff}
\footnotesize\setlength{\tabcolsep}{4pt}
\begin{tabular}{@{}lcccccc@{}}
\toprule
 & \multicolumn{3}{c}{50\% level} & \multicolumn{3}{c}{80\% level} \\
\cmidrule(lr){2-4}\cmidrule(lr){5-7}
Game & Nature & FSQ-J & FSQ-F & Nature & FSQ-J & FSQ-F \\
\midrule
\tabinput{tables/sample_eff_rows.tex}
\bottomrule
\end{tabular}
\end{table}

\paragraph{FSQ-VAE, Frozen.} With a frozen FSQ encoder the agent still learned something in
every game, but much less (median HNS 0.31). It only came close to the NatureCNN on Pong and
MsPacman. It struggled most where the important object is small and fast (Breakout,
DemonAttack, BeamRider) or where the screen scrolls (RoadRunner, Riverraid). Even with the
motion-aware loss, reconstruction still spends most of its capacity on large areas.

\paragraph{V-JEPA.} Both V-JEPA agents learned very little (median HNS 0.05 frozen and 0.04
fine-tuned), even though the frozen critic predicted the returns quite well (explained
variance 0.56--0.97). Fine-tuning was also unstable: on Pong the agent reached $+8.4$ after 7\,M
steps and then fell back to $-15$. My first guess was that the 9\,216-dimensional feature was
too large, so I repeated all 22 runs with the tokens averaged into a
256-dimensional feature and separate encoders for actor and critic, and repeated the frozen
runs once more with a five times larger entropy bonus (\cref{tab:jepa-abl}). Neither change
helped: the pooled frozen agent beat the flattened one on only one game, the pooled fine-tuned
policy became almost deterministic on ten games, and the larger entropy bonus kept the policy
random but stopped the critic from learning.

\begin{table}[t]
\centering
\caption{V-JEPA feature ablation (11 seen games): median HNS of the training return and of
the final policy on 30 test episodes, end-of-training critic explained variance and policy
entropy, and games on which the pooled variant beats the flattened one.}
\label{tab:jepa-abl}
\footnotesize\setlength{\tabcolsep}{3pt}
\begin{tabular}{@{}lcccccccc@{}}
\toprule
 & & & \multicolumn{2}{c}{median HNS} & expl. & & & \\
\cmidrule(lr){4-5}
Feature & mode & ent.\ coef. & train & test & var. & entropy & FPS & wins \\
\midrule
\tabinput{tables/jepa_ablation_rows.tex}
\bottomrule
\end{tabular}
\end{table}

\paragraph{Frozen vs.\ Fine-Tuned and Unseen Games.} For the FSQ encoder, fine-tuning beat
freezing on 13 of 15 games, sometimes by a large factor (Breakout $12\times$, DemonAttack
$7\times$). It also helped on the four held-out games (median HNS 0.45 vs.\ 0.21; Alien 2\,071,
IceHockey $-3.4$). Since I did not train a NatureCNN on these games, I cannot say that the
pre-training actually helped there, only that it did not stop the agent from learning.

\subsection{World-Model Diagnostics}
\label{sec:jepa-diag}

A high similarity between predicted and real latents can be misleading: a collapsed encoder
makes all latents similar, and since $s_{t+1}$ shares three frames with $s_t$, simply copying
$z_t$ already scores well. I therefore tested the V-JEPA on 64 sequences of five steps
and compared it with a copy baseline, an untrained model, and deliberately wrong actions
(\cref{tab:jepa}).

\begin{table}[t]
\centering
\caption{V-JEPA diagnostics: seen games (validation split, true game embedding) and held-out
games (UNKNOWN embedding); chance or untrained values in brackets.}
\label{tab:jepa}
\footnotesize\setlength{\tabcolsep}{4pt}
\begin{tabular}{@{}lcc@{}}
\toprule
Diagnostic & 11 seen games & 4 held-out games \\
\midrule
Cos.\ sim.\ predictor / identity, $k=1$ & 0.987 / 0.974 & 0.939 / 0.971 \\
Cos.\ sim.\ predictor / identity, $k=5$ & 0.961 / 0.915 & 0.806 / 0.912 \\
Action sensitivity: sim.\ true $-$ wrong action & +0.0057 & +0.0001 \\
Attention on action / game token (uniform 0.026) & 0.029 / 0.046 & 0.026 / 0.055 \\
Inverse-dynamics accuracy & 55.9\% (12.9\%) & 1.8\% (5.6\%) \\
Effective rank / 256 & 52.5 (2.5) & 29.9 (1.5) \\
\bottomrule
\end{tabular}
\end{table}

On the seen games the latent space did not collapse (effective rank 52.5 compared with 2.5 for
an untrained model), and the predictor was better than copying at every step. The gap even grew
from $+0.013$ after one step to $+0.046$ after five steps, when the input and target no longer
share any frames, so the model did learn real dynamics. However, the action had very little
influence: giving the predictor a wrong action changed its output by only 0.006, and the
action token received hardly more attention than an average token. The game embedding was
also unused, because the encoder recognised the game from the pixels. On the
held-out games the predictor was worse than copying and the inverse head was below chance, so
the model had learned the eleven training games and not Atari in general.

\subsection{Why the V-JEPA Failed with PPO}
\label{sec:jepa-fail}

Judged by its own training goals, the V-JEPA worked well, but PPO could not use its features.
\Cref{tab:jepa-fail} summarises the evidence I collected, and \cref{fig:pca} shows the main
problem visually.

\begin{table}[!ht]
\centering
\caption{Evidence for why the V-JEPA encoder did not support PPO (from
\cref{sec:pretrain-res,sec:jepa-diag} and \cref{tab:jepa-abl}).}
\label{tab:jepa-fail}
\scriptsize\setlength{\tabcolsep}{3pt}
\newcolumntype{Y}[1]{>{\raggedright\arraybackslash}p{#1}}
\begin{tabular}{@{}Y{2.6cm}Y{5.0cm}Y{4.2cm}@{}}
\toprule
Explanation & Evidence & Consequence for PPO \\
\midrule
Dynamics ignore the action & true vs.\ wrong action: $\Delta$cos.\ 0.006; action-token attention
  0.029 (uniform 0.026); \cref{fig:pca} & latent need not encode what the agent controls \\
Latents change slowly & $\cos(z_t,\bar z_{t+1})=0.974$; 0.915 five steps later;
  $14{\times}14$ patches & small, decisive sprite motion barely moves the feature \\
Organised by game & t-SNE clusters by game; game embedding unused & capacity spent on
  information that is constant within a run \\
Specialised to pre-training games & held-out loss $3\times$ seen; held-out inverse accuracy
  below chance & tuned to expert states of eleven games \\
State, not action, information & critic expl.\ var.\ 0.56--0.97 while policy entropy stays
  high (1.53 vs.\ 1.21) & value readable, best action not \\
Little action variety in the data & expert's top action: 45\% (Qbert), 67\% (RoadRunner) of
  steps (\cref{tab:games}) & few examples of what \emph{other} actions would change \\
Not size or exploration & pooled 256-d feature, entropy bonus 0.05: no gain & rules out the
  simplest remedies \\
\bottomrule
\end{tabular}
\end{table}

\begin{figure}[!ht]
\centering
\includegraphics[width=0.36\textwidth]{figures/pca_action_sensitivity.png}
\caption{PCA of V-JEPA latents for five states: context $z_t$ (square), true next latent
(star) and the prediction for every valid action (dots). The predictions for all actions form
one tight cluster near $z_t$: the predictor hardly distinguishes between actions.}
\label{fig:pca}
\end{figure}

I think the main reason is a mismatch between what the JEPA is trained to do and what a policy
needs. A JEPA predicts what is \emph{predictable} about the next frame, and it can ignore
details that do not change the latent error much~\cite{lecun2022path,assran2023ijepa}. In Atari,
these details, such as the exact position of the ball or of a bullet, are exactly what the agent
needs in order to act. A representation that is good for control should capture what the
agent's actions change~\cite{gelada2019deepmdp}, and this is what my predictor
modelled worst (\cref{fig:pca}). The inverse loss was meant to fix this, but it only trains the
encoder (\cref{tab:grad}), and the expert data contains few examples of alternative actions.
The overlapping frames and the slowly moving EMA target also push the model towards latents
that change very little from one step to the next. Such features are good enough to estimate
how good a state is (the critic's explained variance was high), but too smooth to separate two
states whose best actions differ by a few pixels. Finally,
fine-tuning a 12.5\,M-parameter transformer with PPO's learning rate was unstable, which agrees
with reports that pre-trained visual features transfer poorly to
control~\cite{hansen2022pretrained,parisi2022pvr}.

\FloatBarrier
% =========================================================================
\section{Conclusion}
\label{sec:conclusion}

In this project, I investigated whether one pre-trained encoder can serve as a universal
feature extractor for Atari games. I collected expert data from 15 games, trained VAE, VQ-VAE,
FSQ-VAE and V-JEPA encoders, and used them with PPO for 10\,M steps, either frozen or
fine-tuned. My answer is: partly.

The FSQ-VAE, pre-trained without any rewards on eleven games, turned out to be a good starting
point (\textbf{RQ1}). Fine-tuned with PPO, it was on par with a NatureCNN trained from scratch,
with a higher median score, better final scores on four games and clearly faster learning on
four. Kept frozen, it reached less than a third of that median score, so fine-tuning was
clearly the better choice (\textbf{RQ2}). A frozen encoder would have to keep every small detail that
matters in every game, which a reward-free objective has no reason to do. Neither
model generalised to unseen games without fine-tuning (\textbf{RQ3}). Eleven games are not many,
and the held-out games look quite different (Asteroids is 95\% black). Freezing also saved
little computation, because the game simulation dominates: NatureCNN ran at 976 steps per
second and the frozen V-JEPA at 854 (\textbf{RQ4}). The FSQ-VAE is also easier to inspect,
because its decoder shows what the features contain.

The V-JEPA learned a healthy latent space and real dynamics, but its objective rewarded
predicting what usually happens in a game rather than what the agent's actions change, so its
features were not useful for PPO, whether flattened or pooled. An important lesson for me was that
every bug and shortcut in the V-JEPA pipeline still produced a nice-looking loss curve, so
checks such as the copy baseline and wrong actions were necessary.

This work has some limitations. All results come from a single seed, the FSQ agents and the
other agents used two different PPO implementations, there is no NatureCNN baseline on the
held-out games, and PPO was not re-tuned for the pre-trained encoders. In future work, I would
first repeat the main comparison with three seeds on one implementation. For the V-JEPA,
smaller patches and a training signal in which the action really matters (for example
multi-step predictions or a reward head, as in TD-MPC2) look most promising, and the features
should first be checked by predicting game-state variables such as object
positions. Using the FSQ grid inside a world
model~\cite{hafner2021dreamerv2,micheli2023iris} and pre-training on more games are further
directions. A frozen, truly universal Atari feature extractor is still an open problem.

\paragraph{Acknowledgements.}
I would like to thank Oleh Shkalikov for proposing and supervising this project, and the
Centre for Information Services and High Performance Computing (ZIH) at TU Dresden for the GPU
time on the \emph{Capella} cluster.

\bibliographystyle{splncs04}
\bibliography{references}

\end{document}
